\begin{figure}[H]
\centering
\begin{tikzpicture}[x=.67cm,y=.67cm,font=\footnotesize,>=Latex]
 \draw[step=1,black!12,line width=.3pt] (1,1) grid (9,9);
 \draw[->,black!60] (.7,1)--(9.5,1) node[right] {$i$};
 \draw[->,black!60] (1,.7)--(1,9.5) node[above] {$j$};
 \foreach \k in {1,...,9}{
  \node[below=2pt] at (\k,1) {$\k$};
  \node[left=2pt] at (1,\k) {$\k$};
 }
 \draw[black!65] (1,9)--(9,1);
 \draw[<->,line width=.9pt] (2,8)--(8,2);
 \foreach \x/\y in {2/8,5/5,8/2}{
  \filldraw[fill=white,line width=.8pt] (\x,\y) circle (3pt);
 }
 \fill (5,5) circle (1.5pt);
 \node[above right=3pt,fill=white,inner sep=1pt] at (2,8) {$(2,8)$};
 \node[above right=3pt,fill=white,inner sep=1pt] at (5,5) {$(5,5)$};
 \node[above right=3pt,fill=white,inner sep=1pt] at (8,2) {$(8,2)$};
 \node[anchor=west,align=left] at (10.5,7.6)
   {$i+j=10$\\$R^+=1,\quad q^+=1$};
 \node[anchor=west,align=left] at (10.5,5.4)
   {centro: $ij=25$\\$R^\times=7,\quad q^\times=2$};
 \node[anchor=west,align=left] at (10.5,3.1)
   {vacancias: $ij=16$\\$R^\times=7,\quad q^\times=1$};
\end{tikzpicture}
\caption{Centro electrónico y vacancias conjugadas en la carta APP.
La diagonal es la fibra aditiva \(\mathcal F_{10}\); el centro
es el punto fijo de \((i,j)\mapsto(10-i,10-j)\). Los dos
desplazamientos de módulo \(3\) conservan la suma y el residuo
producto y reducen en una unidad el cociente multiplicativo. Las
flechas representan las deformaciones sobre la fibra.}
\label{el-centro:figura}
\end{figure}
